Mispronunciation detection and diagnosis~(MDD) is a significant part in modern computer-aided language learning~(CALL) systems.
%Phoneme-level pronunciation assessment is key to helping L2 learners improve their pronunciation.
%However, most systems are based on a form of goodness of pronunciation (GOP) which requires pre-segmentation of speech into phonetic units.
Most systems implementing phoneme-level MDD through goodness of pronunciation (GOP), however, rely on pre-segmentation of speech into phonetic units.
This limits the accuracy of these methods and the possibility to use modern CTC-based acoustic models for their evaluation.
In this study, we first propose self-alignment GOP (GOP-SA) that enables the use of CTC-trained ASR models for MDD.
Next, we define a more general segmentation-free method that takes all possible segmentations of the canonical transcription into account (GOP-SF).
We give a theoretical account of our definition of GOP-SF, an implementation that solves potential numerical issues as well as a proper normalization which allows the use of acoustic models with different peakiness over time.
%makes the method applicable with acoustic models with different peakiness over time.
We provide extensive experimental results on the CMU~Kids and speechocean762 datasets comparing the different definitions of our methods, estimating the dependency of GOP-SF on the peakiness of the acoustic models and on the amount of context around the target phoneme.
Finally, we compare our methods
with recent studies over the speechocean762 data
showing that the feature vectors derived from the proposed method achieve state-of-the-art results on phoneme-level pronunciation assessment. 

